\documentclass[../readings.tex]{subfiles}

\begin{document}
\subsection{Root Finding}
\label{subsec:root-finding}

Iterative methods for solving $f(x) = 0$ for continuous $f : \fR \to \fR$.

\subsubsection{Bracketing Methods}

Guaranteed convergence to a root within a known interval.

\addThm{Bisection Method}{%
Given $f \in C[a, b]$ with $f(a)f(b) < 0$, a root $x^* \in (a, b)$ exists.
The method repeatedly halves the interval:
\begin{enumerate}
    \item $m = (a + b)/2$
    \item If $f(a)f(m) < 0$, new interval is $[a, m]$, else $[m, b]$.
\end{enumerate}
\textbf{Error Bound:}\enspace After $n$ steps, $|x_n - x^*| \leq (b - a)/2^n$.
}

\addNote{%
\textbf{Stability:}\enspace Bisection is the most robust method; it always converges if a bracket exists.
\textbf{Convergence:}\enspace **Linear** with rate $1/2$. Gains exactly 1 bit of accuracy per step.
}

\subsubsection{Open Methods}

Faster local convergence, but sensitive to the initial guess $x_0$ and may diverge.

\addDef{Newton's Method}{%
Linearizes $f$ at the current iterate:
\[
    x_{k+1} = x_k - \frac{f(x_k)}{f'(x_k)}.
\]
}

\addThm{Quadratic Convergence}{%
If $f''(x)$ is continuous near a simple root $x^*$ ($f'(x^*) \neq 0$), Newton's method converges **quadratically**:
\[
    |x_{k+1} - x^*| \approx \frac{|f''(x^*)|}{2|f'(x^*)|} |x_k - x^*|^2.
\]
\textbf{Insight:}\enspace The number of correct digits roughly doubles each step.
}

\addNote{%
\textbf{Failure Modes:}\enspace
1. $x_0$ too far from $x^*$.
2. $f'(x_k) \approx 0$ (zero derivative).
3. Multiple roots ($f'(x^*) = 0$): convergence degrades to linear with rate $1/2$.
}

\addDef{Secant Method}{%
Approximates $f'$ via finite differences of the last two iterates:
\[
    x_{k+1} = x_k - f(x_k)\frac{x_k - x_{k-1}}{f(x_k) - f(x_{k-1})}.
\]
Requires two initial points ($x_0, x_1$).
}

\addThm{Superlinear Convergence}{%
The secant method converges with order $\phi = (1 + \sqrt{5})/2 \approx 1.618$ (the golden ratio).
\textbf{Efficiency:}\enspace Faster than bisection, slower than Newton, but requires only 1 function evaluation per step (Newton requires 2: $f$ and $f'$).
}

\subsubsection{Brent's Method}

The standard robust choice for root-finding in 1D.

\addThm{The Hybrid Approach}{%
Brent's method combines:
\begin{enumerate}
    \item \textbf{Bisection:}\enspace For guaranteed global convergence.
    \item \textbf{Secant:}\enspace For fast local convergence.
    \item \textbf{IQI:}\enspace Inverse Quadratic Interpolation using 3 points.
\end{enumerate}
\textbf{Rule:}\enspace It uses fast open steps by default, falling back to bisection only if the step is outside the bracket or not improving quickly enough.
}

\addNote{%
\textbf{Implementation:}\enspace Used in \texttt{scipy.optimize.brentq}.
}

\subsubsection{Exercises}

\addExcr{%
\begin{enumerate}
    \item Bisection on $f(x) = x^3 - 2$ on $[1, 2]$ for 4 steps. Find the guaranteed error.
    \item Newton on $f(x) = x^2 - 2$ from $x_0 = 1$. Verify the doubling of digits.
    \item Prove that for $f(x) = x^2$, Newton converges to $0$ with linear rate $1/2$.
    \item Find the root of $e^x = 3x$ near $x=1$ via Secant method. Compare evaluations vs.\ Newton.
\end{enumerate}
}

\end{document}
